\documentclass[11pt]{article}


% Useful packages
\usepackage[utf8]{inputenc}  % para acentos y ñ
\usepackage[english,spanish]{babel}

\usepackage[letterpaper,top=2cm,bottom=2cm,left=2cm,right=2cm,marginparwidth=1.75cm]{geometry}

\usepackage{amsmath}
\usepackage{amsfonts}
\usepackage{amssymb}

\usepackage{listings} % poner codigo
\lstset{
	inputencoding=utf8,
	extendedchars=true,
	literate={á}{{\'a}}1 {é}{{\'e}}1 {í}{{\'i}}1
	{ó}{{\'o}}1 {ú}{{\'u}}1 {ñ}{{\~n}}1
	{Á}{{\'A}}1 {É}{{\'E}}1 {Í}{{\'I}}1
	{Ó}{{\'O}}1 {Ú}{{\'U}}1 {Ñ}{{\~N}}1,
}


\usepackage{float}

\usepackage{graphicx}

%%%%%%%%%%%%% Referencias
\usepackage[colorlinks=true, allcolors=black]{hyperref}
\usepackage{cleveref}

\graphicspath{ {images/} }

\usepackage{geometry}

\usepackage{fancyhdr}
%\setlength{\headheight}{14pt}
\pagestyle{fancy} 

\usepackage{multicol}
\setlength{\columnsep}{0,5cm}

\usepackage{lipsum}

\usepackage[usenames]{color}

%tablas
\usepackage{tabularx} %ancho de tablas
\usepackage{longtable}
\usepackage{ltablex}
\keepXColumns

\usepackage{caption} % títulos de las tablas/figuras
\captionsetup{font=small} % letra de todos los caption en pequeña
\captionsetup[table]{name=Tabla}
\usepackage{subcaption} % este y el anteriori para poner imagenes juntas

%Bibliografia
\usepackage{hyperref}
\usepackage{cite}
%\usepackage{balance}%evita que se partan cosas

%%%%%%%%%%%%%%% NUEVOS COMANDOS
\newcommand{\comillas}[1]{"#1"} %comillas dobles
\newcommand{\mult}[0]{\!\cdot\!} %comillas dobles


%%%%%%%%%%%%%%%%%%%%%%%%%%%%%%%%%%%%%%%%%%%%%%%%%%%%%%%%%%%
%%%%%%%%%%%%%%%%
%%%%%%%%%%%%%%%%%%%%%%%%%%%%%%%%%%%%%%%%%%%%%%%%%%%%%%%%%%%%%%%%%%%%%%%%%%%
\begin{document}
	
	\begin{center}
		\vspace*{3,5cm}
		\textbf{\Huge Estudio del movimiento browniano mediante el uso de métodos de Monte Carlo}
		
		\bigskip\bigskip\bigskip
		\LARGE{Mario Carracedo Andrés}
		
		\bigskip\bigskip\bigskip\bigskip
		\includegraphics[height=8cm]{imagenes/movimiento_browniano} % Asegúrate de tener esta imagen
	\end{center}
	
	\newpage
	\pagestyle{fancy}
	\fancyhead[l]{Mario Carracedo Andrés}
	\fancyfoot[c]{\thepage}
	
	\section{Qué es el movimiento browniano}
		El movimiento browniano es un proceso estocástico de tiempo continuo que describe el movimiento aleatorio de partículas suspendidas en un fluido. Matemáticamente, se define como un proceso de Wiener $W_t$ que cumple tres propiedades fundamentales:
		\begin{itemize}
			\item $W_0 = 0$ (comienza en el origen).
			\item Trayectorias continuas en casi todo punto.
			\item \textbf{Incrementos independientes y estacionarios:} Para cualquier intervalo de tiempo $\Delta t > 0$, el cambio en la posición sigue una distribución normal: 
			\begin{equation}
				W_{t+\Delta t} - W_t \sim \mathcal{N}(0, \Delta t)
			\end{equation}
			Esto implica que la distancia recorrida es un número aleatorio con valor esperado nulo y varianza igual al tiempo transcurrido. Físicamente, esto refleja que en intervalos de tiempo menores, la partícula tiene menos capacidad de alejarse de su posición actual.
		\end{itemize}
		
		De esta manera se puede decir que $W_t$ representa el movimiento browniano estándar:
		\begin{equation*}
			dW_t = Z_t \sqrt{dt}
		\end{equation*}
		
		Donde $Z \sim \mathcal{N}(0,1)$
		
	\section{Notación y Operadores Estocásticos}
		Para el desarrollo del análisis matemático de la simulación, utilizaremos la siguiente notación estándar para describir las propiedades estadísticas del proceso:
		
		\begin{itemize}
			\item \textbf{Esperanza Matemática o Valor Esperado ($E[X]$):} Representa el promedio ponderado de todos los resultados posibles. En nuestro modelo, define la trayectoria media o \textit{drift}:
			\begin{equation}
				E[X_t] = \mu \mult t
			\end{equation}
			
			\item \textbf{Varianza ($Var(X)$):} Mide la dispersión de los resultados respecto a la esperanza. En el movimiento browniano, la varianza es lineal respecto al tiempo \hyperref[oksendal2003stochastic]{\cite{oksendal2003stochastic}}:
			
			\begin{equation}
				Var(X_t) = \sigma^2 \mult t
			\end{equation}
			
			\item \textbf{Desviación Típica/Estandar ($\sigma_x$):} Es la raíz cuadrada de la varianza. Determina el ancho de la campana de Gauss en cada instante $t$:
			\begin{equation}
				\sigma_x = \sqrt{Var(X_t)} = \sigma \mult \sqrt{t}
			\end{equation}
			
			\item \textbf{Distribución Normal ($\mathcal{N}(\mu, \sigma^2)$):} Denotamos que una variable sigue una distribución gaussiana (campana de Gauss) con una media y una varianza específicas.
			
			\item \textbf{Distribución del Proceso ($X_t \sim \mathcal{N}(\mu t, \sigma^2 t)$):} Denotamos que la posición de la partícula en un instante $t$ sigue una distribución gaussiana. A diferencia de una variable aleatoria estática, los parámetros de esta distribución evolucionan con el tiempo: la media se desplaza según $\mu t$ y la varianza se expande según $\sigma^2 t$ \cite{oksendal2003stochastic}.
		\end{itemize}
		
		Esta distinción es crucial para la simulación numérica, ya que mientras la tendencia del sistema crece de forma lineal con $t$, la incertidumbre o \comillas{ruido} crece de forma sublineal con $\sqrt{t}$.
		
	\section{Random Walk (Caminata aleatoria)}	
		La aproximación discreta del movimiento browniano es el \textit{random walk}\cite{shreve2005stochastic}. En este modelo, una partícula realiza saltos aleatorios en intervalos de tiempo discretos $\Delta t$. Si definimos el salto como $\Delta X = \mu\Delta t + \sigma \sqrt{\Delta t} \mult Z$, donde $Z $ es una variable aleatoria de una distribución normal estándar, el límite cuando $\Delta t \to 0$ converge al movimiento browniano continuo. Este es el principio que utilizamos en nuestra simulación numérica mediante el esquema de Euler-Maruyama.
		
	\section{Conexión con la difusión}
		Existe una equivalencia profunda entre el comportamiento microscópico aleatorio y la ley macroscópica de la difusión. Si consideramos una densidad de partículas $\phi(x,t)$, su evolución sigue la ecuación de calor:
		\begin{equation}
			\frac{\partial \phi}{\partial t} = D \frac{\partial^2 \phi}{\partial x^2}
		\end{equation}
		Donde el coeficiente de difusión $D$ está relacionado con la varianza de la siguiente manera $Var(X_t) = 2D*t$ \hyperref[einstein1906theory]{\cite{einstein1906theory}}, que si lo conectamos con la volatilidad de nuestra simulación,$D = \frac{\sigma^2}{2}$. Esto implica que, aunque una partícula individual es impredecible, el promedio de una población sigue un patrón determinista.
		
	\section{Ecuación diferencial estocástica}
		Para modelar el movimiento con una tendencia (drift), utilizamos la ecuación de Langevin:
		\begin{equation}
			dX_t = \mu dt + \sigma dW_t
		\end{equation}
		Donde:
		\begin{itemize}
			\item $\mu dt$: Representa la componente determinista o \textit{drift}. Es el valor esperado del movimiento por unidad de tiempo.
			\item $\sigma dW_t$: Representa la componente estocástica o difusión.
		\end{itemize}
		
		Integrando esta ecuación, obtenemos que la posición en un tiempo $t$ sigue una distribución normal:
		
		\begin{equation}
			X_t \sim \mathcal{N}(\mu t, \sigma^2 t)
		\end{equation}
		
		Esto justifica que en nuestra simulación hayamos representado el valor esperado como una recta $\mu t$ y la incertidumbre (desviación típica) como una curva proporcional a $\sqrt{t}$.
		
	\section{Implementación Numérica: El Esquema de Euler-Maruyama}
	Para resolver las ecuaciones diferenciales estocásticas en un ordenador, utilizamos el método de Euler-Maruyama. Este discretiza el tiempo en pasos $\Delta t$ y aproxima el incremento del proceso mediante:
	\begin{equation}
		\Delta X = \text{Tendencia}(\Delta t) + \text{Difusión}(\sqrt{\Delta t}, Z)
	\end{equation}
	Donde recordemos que $Z \sim \mathcal{N}(0,1)$. A continuación, se presentan los tres modelos implementados en este proyecto, cada uno con su script de Python independiente. Se hablará algo mas detalladamente sobre el movimiento btowniano geométrico en la subsección \ref{MB geom}
	
	\section{Casos de Estudio y Scripts}
	
	\subsection{Caso 1: Movimiento Browniano Estándar (\texttt{movimiento\_browniano.py})}
	Es el bloque constructivo básico (Proceso de Wiener) donde no existe tendencia ($\mu=0$) y la volatilidad es unitaria ($\sigma=1$). Describe un ruido puramente aleatorio e isotrópico.
	\begin{equation}
		dX_t = dW_t \implies X_t \sim \mathcal{N}(0, t)
	\end{equation}
	En este modelo, la varianza crece exactamente igual al tiempo transcurrido ($Var = t$), y la posición más probable es siempre el origen, aunque la dispersión aumenta con $\sqrt{t}$.
	
	\subsection{Caso 2: Movimiento Browniano Aritmético (\texttt{movimiento\_browniano.py})}
	Este script implementa la ecuación de Langevin, añadiendo una componente determinista (\textit{drift}) que empuja la partícula en una dirección preferente.
	\begin{equation}
		dX_t = \mu dt + \sigma dW_t \implies X_t \sim \mathcal{N}(\mu t, \sigma^2 t)
	\end{equation}
	Aquí, el valor esperado evoluciona como una recta $E[X_t] = \mu t$, mientras que la incertidumbre se expande según las bandas de confianza $\pm \sigma \sqrt{t}$. Es el modelo físico por excelencia para la difusión con corriente o transporte.
	
	\subsection{Caso 3: Movimiento Browniano Geométrico (\texttt{movimiento\_browniano\_geometrico.py})}
	\label{MB geom}
	
	Para modelar activos financieros, utilizamos el Movimiento Browniano Geométrico (GBM). A diferencia de los anteriores, el cambio en el precio es proporcional al valor actual del activo ($S_t$), lo que garantiza que los precios sean siempre positivos ($S_t > 0$).
	\begin{equation}
		dS_t = \mu S_t dt + \sigma S_t dW_t
	\end{equation}
	Usando el Lema de Itô, la solución exacta para el Movimiento Browniano es:
	\begin{equation}
		\boxed{S_t = S_0 e^{ \left(\mu - \frac{\sigma^2}{2}\right)t + \sigma W_t}}
	\end{equation}
	Este script simula el comportamiento de acciones en el mercado, donde la variable sigue una distribución \textbf{log-normal}; es decir, si en el caso aritmético era $X_t \sim \mathcal{N}(\mu t,\sigma^2t)$, ahora es $\ln(S_t) \sim \mathcal{N}\left( \ln(S_0) + \left( \mu - \dfrac{\sigma^2}{2} \right)t  , \sigma^2 t \right)$. Es la base matemática del modelo de Black-Scholes.
	
	De esta manera, ahora $Var(S_t) \neq \sigma^2 \mult t$ y $E[S_t] \neq \mu \mult t$, pues ahora hay que tener en cuenta que la distribución es normal en el logaritmo, de esta manera 
	\begin{equation*}
		\boxed{E[S_t] = S_0e^{ \mu t}}
	\end{equation*}
	\begin{gather*}
		\boxed{Var(S_t) = S_0^2 e^{2\mu t}(e^{\sigma^2t} - 1) \implies  \sigma_{S_t} = \sqrt{Var(S_t)} = S_0 e^{\mu t}\sqrt{e^{\sigma^2t} - 1)}
		}
	\end{gather*}
	
	Si representamos la dispersión absoluta
	
	\begin{gather}
		\boxed{S_{t,up/down} = E[S_t] \pm 1.96 S_0 e^{\mu t}\sqrt{e^{\sigma^2t} - 1)}}
	\end{gather}
	
	A mayores, ese representará también la mediana \comillas{Md} y el intervalo de confianza Log-Nromal (zona en la que el 95\% de caerán ). Esto es debido a que por el comportamiento exponencial, valores que exploten aportarán mucho mas que muchos valores que no lo hagan en, afectando al valor esperado. Si tenemos en cuenta que $Md(W_t) = 0$ (pues la mediana de una distribución normal es el propio valor esperado, y en la distribución normal, este es 0)
	
	\begin{gather*}
		Md(S_t) = Md\left( S_0 e^{ \left(\mu - \frac{\sigma^2}{2}\right)t + \sigma W_t}\right)  = S_0 e^{ \left(\mu - \frac{\sigma^2}{2}\right)t}Md(e^{\sigma W_t}) = S_0 e^{ \left(\mu - \frac{\sigma^2}{2}\right)t}e^{Md(\sigma W_t)} = 
	\end{gather*}
	\begin{equation}
		\boxed{Md(S_t) = S_0 e^{ \left(\mu - \frac{\sigma^2}{2}\right)t}}
	\end{equation}
	
	Veamos la zona de confianza \textbf{Log-Normal}. Como es en el mundo de los logaritmos donde existe la aditividad y se cumple lo que vimos para el caso del browniano aritmetico, $\sigma_{\ln(s_t)} = \sigma\sqrt{t}$, entonces para una banda del 95\% de confianza, se toma margenes $\pm 1.96\sigma_{\ln(s_t)}$
	
	\begin{gather}
		\boxed{S_{t,up/down} = Md(S_t)e^{\pm \sigma_{\ln(s_t)}} = S_0 e^{ \left(\mu - \frac{\sigma^2}{2}\right)t \pm 1.96\sigma \sqrt{t}}}
	\end{gather}
	
	\renewcommand{\refname}{Bibliografía}
	\bibliographystyle{plain}  % O puedes usar otro estilo: ieeetr, apa, etc.
	\bibliography{Bibliografia}
\end{document}
